\section{Related Work}
\label{sec:related-work}

\paragraph{Looped and parameter-shared Transformers.}
Using one block repeatedly has been explored as recurrent depth in
Universal Transformers~\citep{dehghani2018universal} and as cross-layer
parameter sharing in ALBERT~\citep{lan2019albert}. More recent looped
models show strong behavior on iterative algorithm learning and
multi-step in-context procedures~\citep{yang2023looped,gatmiry2024multistep},
length generalization~\citep{fan2024looped}, and latent-reasoning style
test-time compute scaling~\citep{saunshi2025latent,zhu2025ouro}. Related
parameter-sharing formulations also appear in looped neural networks~\citep{ng2024loopnn}. Our work addresses the complementary question of how to parameterize the residual scaling so that training remains stable and hyperparameters transfer.

\paragraph{Stability in deep residual stacks and deep Transformers.}
Large-depth training has a long line of stabilization techniques,
including residual reparameterizations and initialization rules such as
Fixup and ReZero~\citep{zhang2019fixup,bachlechner2020rezero}, and
Transformer-specific stabilizers such as DeepNet/DeepNorm~\citep{wang2022deepnet}. On the theory side, prior analyses of deep
non-shared residual networks characterize when residual scaling keeps
signals controlled in the large-depth limit~\citep{marion2022scalingresnets}. A complementary line studies
generalization for continuous-depth models and their ResNet analogues~\citep{DBLP:conf/nips/ChenRBD18}:
\citet{marion2023generalization} derives a Lipschitz-based bound whose
complexity term depends on differences between successive weight
matrices. These works mostly treat depth as a
stack of distinct layers, whereas in our setting loop steps reuse the
same parameters.

\paragraph{Hyperparameter transfer and parameterization.}
Tensor-program and $\mu$P-style analyses establish transfer principles
across model scales and motivate systematic parameterization choices~\citep{yang2022tensorprogramsV}. Recent depth-transfer analyses in
residual networks study how learning-rate and initialization choices
change with depth under non-shared assumptions~\citep{bordelon2024depthwise,hayou2023commute}. CompleteP and follow-up
work extend this direction for deep Transformers and broader axes of
transfer~\citep{dey2025completep,mlodozeniec2025completed}. Our work is complementary, targeting the \emph{loop axis} and showing that shared weights create cross-step correlations that change the stability threshold and transfer regime.
